% ATTEMPT_STATUS: unresolved
% RESEARCH_GENERATED: 2026-10-01; GPT-6 Astra Ultra; individual mathematical attempt
% TOP_PROBLEM: {"id": 173, "title": "Covering by Residue Classes", "category_id": 1, "category": {"id": 1, "name": "number_theory", "display_name": "Number Theory", "description": "Properties of integers, prime numbers, Diophantine equations.", "slug": "number-theory", "order_index": 1, "created_at": "2026-07-31T15:26:25.670Z"}, "status": "open"}
% FIRST_POSTED: 2026-10-01
% Imported from: unsolvedmath_additional_500.tex; UnsolvedMath ID 173
% !TeX program = lualatex
% Compile twice: lualatex 173.tex
% Self-contained: no external bibliography, images, or \input files.
% Numeric section labels and visible numbers are original UnsolvedMath IDs.
\documentclass[11pt,a4paper]{article}
\usepackage[margin=24mm,headheight=14pt]{geometry}
\usepackage{fontspec}
\setmainfont{Latin Modern Roman}
\setsansfont{Latin Modern Sans}
\setmonofont{Latin Modern Mono}
\usepackage{amsmath,amssymb,mathtools}
\usepackage{unicode-math}
\setmathfont{Latin Modern Math}
\usepackage{microtype}
\usepackage{enumitem,xurl,needspace}
\usepackage[dvipsnames]{xcolor}
\usepackage{hyperref}
\hypersetup{unicode=true,colorlinks=true,linkcolor=MidnightBlue,urlcolor=MidnightBlue,
 pdftitle={UnsolvedMath Problem 173},pdfauthor={Source-annotated compilation},bookmarksnumbered=true}
\usepackage{bookmark,tocloft,titlesec,fancyhdr}
\setlength{\cftsecnumwidth}{4.2em}
\cftsetpnumwidth{2.5em}
\cftsetrmarg{4em}
\titleformat{\section}{\Large\bfseries\raggedright}{\thesection}{0.7em}{}
\pagestyle{fancy}\fancyhf{}
\fancyhead[L]{\small\sffamily UnsolvedMath research attempt}
\fancyhead[R]{\small Original catalogue IDs}
\fancyfoot[C]{\thepage}
\setlength{\parindent}{0pt}
\setlength{\parskip}{5pt plus 1pt}
\setlength{\emergencystretch}{3em}
\setcounter{tocdepth}{1}\setcounter{secnumdepth}{1}
\setlist[itemize]{leftmargin=3em,itemsep=4pt,topsep=4pt,parsep=0pt}
\allowdisplaybreaks
\newcommand{\N}{\mathbb{N}}
\newcommand{\Z}{\mathbb{Z}}
\newcommand{\Q}{\mathbb{Q}}
\newcommand{\R}{\mathbb{R}}
\newcommand{\C}{\mathbb{C}}
\newcommand{\F}{\mathbb{F}}
\newcommand{\A}{\mathbb{A}}
\newcommand{\PP}{\mathbb{P}}
\newcommand{\E}{\mathbb{E}}
\newcommand{\eps}{\varepsilon}
\newcommand{\dd}{\,\mathrm d}
\newcommand{\sref}[2]{\hyperlink{source:#1:#2}{[S#2]}}
\newif\ifshowcatalogue
\showcataloguetrue
\newcommand{\cataloguescope}[1]{\ifshowcatalogue
 \paragraph{Statement recorded in the supplied source.}
 #1\fi}
\begin{document}
% UnsolvedMath ID 173; source catalogue code GREEN-085

\setcounter{section}{172}

\section{Coverage by One Residue Class per Medium Prime}\label{173}

{\small\sffamily UnsolvedMath ID 173 · Number Theory}

\subsection{Definitions and mathematical statement}

\paragraph{Shared notation.}Unless a section says otherwise, $\N=\{1,2,\ldots\}$,
$\Z$ is the ring of integers, $\Q,\R,\C$ are the rational, real, and complex
fields, $[N]=\{1,\ldots,\lfloor N\rfloor\}$, and $\log$ is the natural
logarithm. For real functions of a parameter tending to infinity,
$f\ll g$ and $f=O(g)$ mean $|f|\leq Cg$ eventually for some fixed $C$;
$f\gg g$ reverses this comparison; $f=o(g)$ means $f/g\to0$;
$f\sim g$ means $f/g\to1$; and $f\asymp g$ means both order bounds.
A subscript on $O$ or $\ll$ specifies allowed constant dependence.
The expression $n^{o(1)}$ in an upper bound means growth smaller than
$n^\epsilon$ for each fixed $\epsilon>0$. Local definitions govern any
exception, finite-field convention, or use of the same symbol for another
object. An asymptotic claim is not automatically an effective algorithm.



Let $\mathcal P_N$ be the primes in $[N^{0.51},2N^{0.51})$. For a completely arbitrary choice of residues $a(p)$ define $U_N=\{n\in[N]:n\equiv a(p)\pmod p\text{ for some }p\in\mathcal P_N\}$. The bound asks that $|U_N|\geq N^{1-o(1)}$, uniformly in the residue choices; equivalently each fixed power loss eventually suffices. \sref{173}{1} \sref{173}{2}

\paragraph{Statement recorded in the supplied source.}
Let $N$ be large. For each prime $p$ with $N^{0.51} \leq p < 2N^{0.51}$, pick a residue $a(p) \in \mathbb{Z}/p\mathbb{Z}$. Is $\#\{n \in [N] : n \equiv a(p) \pmod p \text{ for some } p\} \gg N^{1-o(1)}$? \sref{173}{1}

\subsection{Short English statement}

Regardless of how one residue is chosen for each prime near a specified power of N, must their union cover almost a full power of the initial interval?

\subsection{Research attempt}

\paragraph{Provenance and scope.}
This individual attempt was generated by GPT-6 Astra Ultra on 1 October 2026. The method was to derive explicit partial results and test the first proposed extension adversarially. The original statement and sources below are retained. No claim of mathematical novelty or complete solution is made.

\paragraph{Formal target and the chosen attack.}
Put $Q=N^{0.51}$. For every prime $p\in[Q,2Q)$ choose one arbitrary residue $a(p)$, and let $U$ be the union of the resulting progressions inside $[N]$. The question seeks $|U|\gg N^{1-o(1)}$ uniformly over all choices. The primary sources place this among large-modulus sieve questions. We derive the exact first-and-second-moment bound obtainable from pairwise intersections, identify its exponent, and test whether higher incidence alone can improve it.

\paragraph{Every pair of progressions intersects at most once.}
Let $\mathcal P$ be the prime set, $m=|\mathcal P|$, and $P_p=\{n\in[N]:n\equiv a(p)\pmod p\}$. For distinct primes $p,q$, two different common points would have difference divisible by $pq$. But $pq\geq Q^2=N^{1.02}>N$, whereas two points of $[N]$ differ by less than $N$. Hence $|P_p\cap P_q|\leq1$, regardless of the residues.

Write $r(n)=\#\{p:n\in P_p\}$ and $I=\sum_nr(n)=\sum_p|P_p|$. Each progression has between $N/p-1$ and $N/p+1$ points, so
\[
 m\left(\frac{N}{2Q}-1\right)\leq I
 \leq m\left(\frac NQ+1\right).
\]
The exact second-moment identity is
\[
 \sum_nr(n)^2=I+2\sum_{p<q}|P_p\cap P_q|
 \leq I+m(m-1).
\]
Cauchy--Schwarz on the support of $r$ now proves the uniform bound
\[
 |U|\geq\frac{I^2}{I+m(m-1)}.
\]
This finite formula contains no distributional hypothesis about the chosen residues.

\paragraph{The exponent it actually gives.}
The classical prime number theorem gives $m\asymp Q/\log Q$ for the dyadic interval. Thus $I\asymp N/\log N$, while $m^2\asymp N^{1.02}/(\log N)^2$ dominates $I$ for large $N$. Substitution yields
\[
 |U|\gg N^{0.98}.
\]
More generally, for moduli around $N^\theta$ with fixed $\theta>1/2$, the same argument gives exponent $2-2\theta$. At $\theta=0.51$ that is a fixed loss of $0.02$, not an $o(1)$ loss. Writing the numerical exponent as nearly one must not obscure this difference in quantifiers.

The bound is stronger than the size of a single progression, which is only order $N^{0.49}$, and remains valid even if all residues are chosen to make one common point lie in every progression. Such concentration shows why counting intersecting pairs can be much less informative than counting the union itself.

\paragraph{Why pairwise-intersection information alone has a barrier.}
There are abstract incidence systems with the same type of parameters and a union of order the square of the line size. For example, in the affine plane over a finite field of order $L$, every line has $L$ points and two distinct lines intersect in at most one point. One may choose between $L$ and $L^2$ lines while their union stays inside the $L^2$-point plane. Taking about $L^{0.51/0.49}$ lines matches, at the exponent level, the progression count $m\approx N^{0.51}$ and line length $L\approx N^{0.49}$. The whole ambient incidence universe has size $L^2\approx N^{0.98}$.

This is not a construction of the required integer residue classes: affine-plane lines do not automatically arise as prime-modulus progressions in an interval. It proves a limitation on an argument that uses only their sizes and the at-most-one pair intersection property. An improvement to near-linear coverage must exploit arithmetic information absent from that abstract model.

\paragraph{A precise missing input.}
The moment formula would yield $|U|\geq N^{1-o(1)}$ if one could prove $\sum_nr(n)^2\leq N^{1+o(1)}$, since $I=N^{1-o(1)}$. The trivial pair bound instead permits $N^{1.02+o(1)}$. But this proposed second-moment estimate is itself false for the fully aligned choice $a(p)=n_0\pmod p$: the single point $n_0$ contributes $m^2=N^{1.02+o(1)}$. Thus even the most obvious strengthened moment target is ruled out. A viable refinement must remove a small exceptional set of very high multiplicity, then control the remaining mass rather than demand an untruncated global second moment.

\paragraph{Verification and status.}
The intersection argument uses distinct primes and the strict product bound, and every floor error is recorded in $I$. The aligned-residue example is allowed by the original arbitrary-choice quantifier. The attempt proves the uniform $N^{0.98}$ bound and disproves a natural untruncated moment strengthening. It leaves unresolved the arithmetic control of the residual incidence mass after high-multiplicity points are separated, which is needed for the requested near-linear lower bound.

\subsection{Sources}

{\small

\begin{itemize}

\item[\hypertarget{source:173:1}{S1}] ULAM.AI, \emph{UnsolvedMath}, supplied \texttt{problems.json}, record 173 (GREEN-085), statement and background. The embedded literature-review date is 2026-08-17; that is the archive’s review date, not a fresh certification. Dataset landing page: \url{https://huggingface.co/datasets/ulamai/UnsolvedMath}.

\item[\hypertarget{source:173:2}{S2}] B. Green, A note on multiplicative functions on progressions to large moduli, Proc. Roy. Soc. Edinburgh Sect. A 148 (2018), 63-77. \url{https://doi.org/10.1017/S0308210517000144}. Bibliographic information imported from the supplied record.

\item[\hypertarget{source:173:3}{S3}] B. Green, 100 Open Problems, Problem 43, version. \url{https://people.maths.ox.ac.uk/greenbj/papers/open-problems.pdf}. The author-hosted document was inspected on 27 September 2026; the archive’s local code is not a problem-number locator in this paper.

\end{itemize}

}

\label{end:173}
\end{document}
